\documentclass{article}
\usepackage{amsmath}
\usepackage{amssymb}
\usepackage{array}
\usepackage{booktabs}
\usepackage{textcomp}
\usepackage{graphicx}
\usepackage{tikz}
\usepackage{graphics}
\usepackage{pdfpages}

\usepackage{amsthm}
\newtheorem{theorem}{Theorem}[section]

\newtheorem{lemma}[theorem]{Lemma}



\newcommand{\coNP}{\textrm{co-}\mathbb{NP}}

\title{cse 242 hw2}
\author{Ian Chi}
\date{\today}
\begin{document}
\maketitle

\section{}
\begin{enumerate}
    \item 
    \begin{enumerate}
        \item \begin{equation}
        \begin{split}
                J(w)&=\Big|Xw-y\Big|_2^2 + \lambda\Big|w\Big|_2^2\\
                J(w)&=(Xw-y)^T(Xw-y) + \lambda w^T w\\
                J(w)&=w^TX^TXw - w^TX^Ty - y^T Xw - y^Ty + \lambda w^T w\\
                J(w)&=w^TX^TXw - 2y^T Xw - y^Ty + \lambda w^T w\\
                \frac{\partial J(w)}{\partial w}&=2X^TXw -2y^TX + 2\lambda w = 0\\
        \end{split}\end{equation}
        
        then we see 
        \begin{equation}
            \begin{split}
                2X^TXw  + 2\lambda w &= 2y^TX\\
                (X^TX  + \lambda I )w &= y^TX\\
            \end{split}
        \end{equation}
        
        then we get the solution in the hw. 

        \item as \(\lambda\to 0\), \(w_{ridge}->(X^TX)^{-1}X^Ty\).  as \(\lambda\to \infty\), \(w_{ridge}->0\). 
        
        the former is when penalty goes to zero and the latter is when error dominates. 

        \item 
    \end{enumerate}
    
\end{enumerate}



\end{document}